\documentclass[11pt,a4paper]{article}
\usepackage[T1]{fontenc}
\usepackage[margin=25mm]{geometry}
\usepackage{amsmath,amssymb,amsthm}
\usepackage{microtype}

\newtheorem{conjecture}{Conjecture}
\newtheorem{lemma}{Lemma}
\newtheorem{theorem}{Theorem}
\pagestyle{empty}
\setlength{\parindent}{0pt}
\setlength{\parskip}{0.45em}

\title{The Equal Halves Conjecture\\[0.3em]
  \large The case of connected Cayley graphs of even order}
\author{Luc Gillibert, GPT 6 Astra}
\date{17/11/2026}

\begin{document}
\maketitle
\vspace{-2em}

All graphs are finite, simple, and undirected. A graph is
\emph{vertex-transitive} if its automorphism group acts transitively on
its vertices. For $A\subseteq V(G)$, the induced subgraph $G[A]$ retains
every edge of $G$ whose endpoints both lie in $A$. A single vertex is
regarded as a connected graph.

\begin{conjecture}[Equal Halves]
Every connected vertex-transitive graph $G$ of even order admits a
partition $V(G)=A\sqcup B$ into two nonempty sets such that $G[A]$ and
$G[B]$ are connected and isomorphic.
\end{conjecture}

The two parts necessarily have equal size. The following lemma proves
the conjectured conclusion whenever a suitable involution exists;
vertex-transitivity is not needed for the lemma.

\begin{lemma}
Let $G$ be a nonempty connected graph admitting an automorphism $\tau$
such that $\tau^2=\mathrm{id}$ and $\tau(v)\ne v$ for every vertex $v$.
Then $V(G)=A\sqcup\tau(A)$ for some $A$ with $G[A]$ connected.
Consequently, the two induced halves are connected and isomorphic.
\end{lemma}

\begin{proof}
Choose a nonempty set $A$ of maximum cardinality subject to $G[A]$
being connected and $A\cap\tau(A)=\varnothing$. Such a set exists,
since any singleton satisfies these conditions.

Suppose $U=A\cup\tau(A)$ is not all of $V(G)$. By connectivity, there
is an edge $ux$ with $u\in U$ and $x\notin U$. The set $U$ is
$\tau$-invariant, so, applying $\tau$ to this edge if necessary, we may
assume $u\in A$ and still $x\notin U$. Both $x$ and $\tau(x)$ lie
outside $U$, and they are distinct. Therefore
\[
  (A\cup\{x\})\cap\tau(A\cup\{x\})=\varnothing.
\]
The edge $ux$ also makes $G[A\cup\{x\}]$ connected, contradicting
maximality. Hence $U=V(G)$. The restriction of $\tau$ is an
isomorphism $G[A]\to G[\tau(A)]$, proving the claim.
\end{proof}

\begin{theorem}
Every connected Cayley graph of even order admits a partition into
two isomorphic connected induced subgraphs.
\end{theorem}

\begin{proof}
Write $G=\operatorname{Cay}(H,S)$, where $H$ is a finite group of even
order, $S=S^{-1}\subseteq H\setminus\{1\}$ generates $H$, and edges
join $h$ to $hs$ for $h\in H$ and $s\in S$.
By Cauchy's theorem, $H$ contains an element $t$ of order two.
Define $\tau(h)=th$. Left multiplication preserves adjacency, since
the edge $\{h,hs\}$ maps to $\{th,(th)s\}$. Thus $\tau$ is a graph
automorphism. Moreover, $\tau^2(h)=t^2h=h$, and $\tau(h)=h$ would imply
$t=1$, a contradiction. The lemma applies and gives the required
partition.
\end{proof}

\end{document}
